\documentclass[conference]{IEEEtran}
\usepackage{amsmath}
\usepackage{booktabs}
\usepackage{pgfplots}
\pgfplotsset{compat=1.17}

\begin{document}

\title{Parallel and Distributed K-Nearest-Neighbor Classification: Comparing POSIX Threads, OpenMP, and MPI}

% Name/email filled in from your git identity; double-check the name spelling.
\author{
\IEEEauthorblockN{Rahat Rizvi Rahman}
\IEEEauthorblockA{Department of Computer Science\\
Virginia Commonwealth University, VCU\\
Richmond, USA\\
rahmanr12@vcu.edu}
}

\maketitle

\begin{abstract}
This report presents a serial baseline and three parallel implementations
of the k-nearest-neighbor (KNN) classifier: a shared-memory version using
POSIX threads, a shared-memory version using OpenMP, and a distributed
version using MPI. All three parallel versions partition the classification
work along the same axis, the set of test instances, so that every worker
reproduces the exact predictions of the serial version regardless of how
many workers are used, which was confirmed exactly on every one of 244
completed runs across two clusters. The implementations were developed and
correctness tested on \texttt{maple.cs.vcu.edu}, then evaluated for runtime
and speedup on \texttt{maple.cs.vcu.edu} and, in place of an inaccessible
\texttt{athena.hprc.vcu.edu} account, on the \texttt{hickory.cs.vcu.edu}
cluster, using three datasets of increasing size. On hickory's dedicated,
uncontended allocation all three implementations scale consistently with
worker count, reaching up to a 5.8$\times$ speedup at 6 workers on the
largest dataset, while identical code on the shared, unscheduled maple
machine shows only modest gains on the larger datasets and measurable
regressions on the smallest one, isolating scheduling and contention,
rather than the parallelization approach itself, as the deciding factor.
\end{abstract}

\begin{IEEEkeywords}
HPDS, KNN, POSIX threads, OpenMP, MPI, parallel computing, speedup
\end{IEEEkeywords}

\section{Introduction}

The k-nearest-neighbor classifier assigns a class to a query instance by
finding the $k$ closest instances to it in a labeled training set and
taking a majority vote of their labels. Its accuracy is competitive with
far more complex classifiers on many problems, but its cost grows as
$O(\mathit{num\_test} \times \mathit{num\_train} \times
\mathit{num\_features})$, since every test instance must be compared
against every training instance across every feature. On datasets with
tens of thousands of instances this cost becomes the practical limit on
how quickly the classifier can be run, which makes KNN a useful vehicle for
comparing parallelization strategies: the algorithm itself is simple
enough that the effect of the parallelization strategy, rather than the
complexity of the algorithm, dominates the results.

This report implements and evaluates three parallel strategies for the
same classification workload. All three are built around the same
partitioning decision, splitting the test instances across workers rather
than the training instances, for a reason argued in Section~\ref{sec:partition}:
it is the only one of the two splits that reproduces the serial version's
predictions exactly, including on ties, without any extra work.

The rest of this report is organized as follows. Section~\ref{sec:methodology}
describes the serial algorithm, the shared partitioning strategy, the three
parallel implementations, and the experimental setup and metrics used to
evaluate them. Section~\ref{sec:results} presents the measured runtimes,
speedups, and accuracy. Section~\ref{sec:analysis} discusses the trends in
those results and compares the three approaches. Section~\ref{sec:conclusion}
concludes.

\section{Methodology}
\label{sec:methodology}

\subsection{Serial KNN Algorithm}

The provided serial implementation classifies each test instance
independently. For a single test instance, it computes the Euclidean
distance to every training instance, maintains a sorted list of the $k$
closest training instances seen so far, and after scanning the full
training set predicts the class that appears most often among those $k$
neighbors. A new training instance is inserted into the candidate list only
when its distance is strictly smaller than an existing candidate's
distance, so among instances at exactly equal distance, the one encountered
earlier in the scan keeps its place. This tie-breaking rule matters for
Section~\ref{sec:partition}.

\subsection{Work Partitioning Strategy}
\label{sec:partition}

All three parallel implementations divide the set of test instances into
contiguous blocks and assign one block to each worker (thread or MPI
process), which then scans the entire training set for each test instance
in its block, exactly as the serial version does. Worker $t$ of $P$ workers
is given the half-open range $[\lfloor tn/P \rfloor, \lfloor (t+1)n/P
\rfloor)$ over $n$ test instances. This splits the test instances into $P$
blocks whose sizes differ by at most one, and it degrades gracefully when
$P$ exceeds $n$: a worker with $t \geq n$ receives an empty range and
simply does no work, which occurs in this evaluation whenever the small
dataset (160 test instances) is run with more than 160 workers.

An alternative partitioning splits the training set instead, and has every
worker scan a shard of the training instances for the same query,
building partial candidate lists that are then merged. This exposes more
parallel work, since the training sets used here are an order of magnitude
larger than the test sets, but it changes the order in which training
instances are compared for a given query. Because the serial algorithm
breaks distance ties by scan order, merging partial candidate lists in a
different order than the original single-threaded scan can change which
neighbor is kept when two training instances are exactly equidistant from
a query, changing the resulting prediction. Splitting on the test set
instead preserves the original scan order inside every worker, so every
implementation in this report reproduces the serial predictions exactly.
This is confirmed empirically in Section~\ref{sec:results}.

Splitting on the test set is also already load balanced without any
runtime scheduling. Every test instance costs exactly
$\mathit{num\_train} \times \mathit{num\_attributes}$ operations to
classify, so an equal-sized static split already assigns equal work to
every worker, and none of the three implementations uses a dynamic
work-stealing or work-queue scheme.

\subsection{POSIX Threads Implementation}

The pthreads version packs every value a worker thread needs, the shared
read-only dataset pointers, the worker's assigned range of test instances,
and a pointer to the shared output array, into a struct, and allocates one
struct instance per thread so that no two threads ever read or write the
same struct while a thread might still be starting up. Each worker writes
predictions only into the slice of the output array given by its own
range, so no lock is needed on the output array, since the ranges assigned
to different threads never overlap.

Each worker also allocates its own private copy of the two scratch
buffers used during classification, the sorted candidate list and the
per-class vote counts, once when the thread starts, and reuses that
private copy across every test instance in its range. The serial version
reuses a single copy of these buffers across all test instances because
only one instance is ever being classified at a time; sharing that same
buffer between multiple threads would let two threads overwrite each
other's in-progress candidate list, so each thread is given its own
independent copy instead. The main thread waits for every worker thread to
finish with \texttt{pthread\_join} before reading the completed
predictions array, which acts as the barrier that guarantees every write
is visible before the result is used.

\subsection{OpenMP Implementation}

The OpenMP version expresses the same partitioning using
\texttt{\#pragma omp parallel} to start a team of threads, followed by a
separate \texttt{\#pragma omp for}, with no \texttt{schedule} clause, to
divide the test-instance loop across that team. The two pragmas are kept separate,
rather than using the combined \texttt{parallel for} form, so that each
thread's private scratch buffers can be allocated once, between the two
pragmas, instead of once per test instance.

Those scratch buffers are declared as ordinary local variables inside the
parallel region rather than passed through OpenMP's \texttt{private} or
\texttt{firstprivate} data-sharing clauses. For a pointer variable, both
clauses only copy the pointer's value, an address, not the buffer that
address points to; every thread would therefore end up pointing at the
same shared buffer, or in the case of \texttt{private}, at an
uninitialized address. Declaring the buffers inside the parallel region
instead relies on the fact that every thread executing that region runs
the declaration itself and so receives its own independent storage,
without needing any explicit data-sharing clause. No \texttt{schedule}
clause is written because, as argued in Section~\ref{sec:partition}, every
test instance costs the same amount of work, so the fixed, contiguous
split every mainstream OpenMP runtime defaults to is already balanced, and
an explicit dynamic schedule would only add scheduling overhead for no
benefit here. The implicit barrier at the end of the \texttt{omp for} construct guarantees
every thread has finished writing its predictions before its scratch
buffers are freed and before the parallel region closes.

\subsection{MPI Implementation}

Unlike the two shared-memory versions, MPI ranks are separate processes
with separate address spaces, so data cannot be shared through a common
pointer. The implementation avoids communicating the dataset itself by
having every rank independently parse both the training and test files
before the timed region begins, so every rank ends up with its own
complete copy of the data without any dataset bytes crossing between
processes. This trades per-rank memory for simplicity and for eliminating
one whole category of communication, at the cost of not scaling to
datasets too large to fit in a single node's memory, which is not a
limitation for the datasets used here.

Every rank computes the same partition table used by the threaded version,
covering every rank's range, not only its own, since the collection step
described next needs the full table. Each rank then classifies its own
range of test instances into a local array sized to exactly its own share
of the work. Because each rank's scratch buffers and local results array
exist only in that rank's own memory, there is no data race to avoid here
the way there is with threads.

The one step with no equivalent in the threaded versions is collecting
every rank's local results back into a single array on rank~0, which alone
computes the confusion matrix. This is done with \texttt{MPI\_Gatherv}
rather than the simpler \texttt{MPI\_Gather}, because \texttt{MPI\_Gather}
requires every rank to contribute the same number of elements, while the
partition table gives blocks that differ by one whenever the test set does
not divide evenly across ranks. \texttt{MPI\_Gatherv} takes a per-rank
count and a per-rank displacement, so each rank's contribution lands at
the correct global offset regardless of block size, and rank~0 recovers
the predictions in their original order. \texttt{MPI\_Gatherv} is also a
synchronization point, so an \texttt{MPI\_Barrier} is placed immediately
before the timer starts, to prevent a rank that happens to finish parsing
its input files earlier than the others from including that head start,
and the resulting idle wait inside \texttt{MPI\_Gatherv}, in its own
measured runtime.

\subsection{Experimental Setup}

Development, debugging, and correctness testing were done on
\texttt{maple.cs.vcu.edu}, a single shared, unscheduled node with 28 logical
CPUs used concurrently by other students, with no reservation mechanism.
Full timing experiments were intended to run on both \texttt{maple.cs.vcu.edu}
and the \texttt{athena.hprc.vcu.edu} Slurm cluster, with the MPI
implementation run twice on Athena, once with every rank forced onto a
single node and once allowed to spread across multiple nodes, to compare
intra-node against inter-node communication cost. This plan could not be
carried out: the assigned Athena account returned access errors for the
full duration of the assignment period despite being reported to the course
staff, so no Athena results appear in this report.

In its place, the \texttt{hickory.cs.vcu.edu} cluster, a separate VCU HPRC
system, was used for the reserved-allocation comparison against maple.
hickory imposes constraints that do not exist on Athena and that limit what
could be measured there. Its Slurm configuration exposes a single node in a
single, GPU-only partition, so there is no second node to place a
multi-node MPI run on; the single-node-versus-multi-node MPI comparison the
assignment specifies for Athena could therefore not be performed on
hickory, and only the single-node MPI case is reported. The partition's
job-submission policy also rejects any job that requests no GPU resource,
so every hickory job requests one unit of the smallest available GPU type
purely to satisfy that policy, without the benchmark ever touching it. The
account's QOS caps a single job at 6 CPUs, so hickory worker counts are
limited to 1, 2, 4, and 6, rather than the full 1--128 range used on maple.

Three datasets of increasing size were used, summarized in
Table~\ref{tab:datasets}. Each was evaluated with $k=3$. On maple, worker
counts of 1, 2, 4, 8, 16, 32, 64, and 128 were attempted; oversubscribing
maple's 28 physical CPUs with dozens to well over a hundred threads or
processes, on a machine already shared with other students, made the
higher worker counts on the two larger datasets take impractically long
(single configurations were observed spending over 95\% of their wall-clock
time waiting for a CPU rather than computing), so the maple results
reported here cover the small and medium datasets at worker counts up to
64 (32 for MPI); the large dataset completed only its serial baseline on
maple before this became impractical. hickory's own worker counts are the
1, 2, 4, 6 range its CPU cap allows. Every configuration was run three
times and the runtimes reported in Section~\ref{sec:results} are the mean
of those three runs.

\begin{table}[htbp]
\caption{Dataset sizes}
\label{tab:datasets}
\centering
\begin{tabular}{lrrr}
\toprule
Dataset & Train instances & Test instances & Attributes \\
\midrule
small  & 1{,}184  & 160   & 7 \\
medium & 14{,}708 & 740   & 11 \\
large  & 61{,}606 & 3{,}436 & 11 \\
\bottomrule
\end{tabular}
\end{table}

\subsection{Performance Metrics}

Two speedup definitions are reported, since they answer different
questions. Absolute speedup compares a parallel run against the serial
program,
\begin{equation}
S_{\text{abs}}(p) = \frac{T_{\text{serial}}}{T_{\text{parallel}}(p)},
\label{eq:abs}
\end{equation}
and reflects the real-world benefit of parallelizing, including whatever
fixed overhead the parallel version pays. Relative speedup instead
compares a parallel implementation against its own single-worker run,
\begin{equation}
S_{\text{rel}}(p) = \frac{T_{\text{parallel}}(1)}{T_{\text{parallel}}(p)},
\label{eq:rel}
\end{equation}
which isolates how well that implementation scales once its own fixed
overhead is factored out. Absolute speedup is used as the headline metric
in Section~\ref{sec:results}, with relative speedup reported alongside it.

The parallel fraction of the code is estimated from the absolute speedup
using the Karp-Flatt metric,
\begin{equation}
e(p) = \frac{1/S_{\text{abs}}(p) - 1/p}{1 - 1/p}, \qquad
f_{\text{parallel}}(p) = 1 - e(p),
\label{eq:karpflatt}
\end{equation}
where $e(p)$ is the estimated serial fraction at $p$ workers. Unlike
Amdahl's law solved for a single fixed serial fraction, Karp-Flatt is
evaluated at every measured value of $p$; an estimated serial fraction that
grows as $p$ increases is itself evidence of parallel overhead, such as
load imbalance or communication cost, that a model with one fixed serial
fraction cannot capture.

\section{Results}
\label{sec:results}

\subsection{Accuracy}

Table~\ref{tab:accuracy} reports the classification accuracy of the serial
implementation on each dataset. Every parallel implementation, at every
worker count tested on either cluster, reproduced these accuracies exactly,
confirmed with an automated check comparing each parallel run's reported
accuracy against the serial run's accuracy for the same dataset; across all
244 completed runs on maple and hickory combined, not one mismatch was
observed. The lower accuracy on the medium dataset is a property of that
dataset's class separability, not a defect in any implementation, since it
is reproduced identically by the serial version and by every parallel run.

\begin{table}[htbp]
\caption{Serial classification accuracy, $k=3$}
\label{tab:accuracy}
\centering
\begin{tabular}{lr}
\toprule
Dataset & Accuracy \\
\midrule
small  & 85.00\% \\
medium & 27.57\% \\
large  & 99.48\% \\
\bottomrule
\end{tabular}
\end{table}

\subsection{Runtime and Speedup on Maple}

Tables~\ref{tab:maple-threaded}--\ref{tab:maple-mpi} report mean runtime,
absolute speedup, and estimated parallel fraction for all three
implementations on maple, for the small and medium datasets; as explained
in Section~\ref{sec:methodology}, the large dataset did not finish beyond
its serial baseline on maple within a practical amount of time and is
omitted from these tables. Figure~\ref{fig:maple-speedup} plots absolute
speedup against worker count for all three implementations on the medium
dataset, the largest one with complete maple coverage.

\begin{table*}[t]
\caption{Pthreads on maple ($k=3$)}
\label{tab:maple-threaded}
\centering
\begin{tabular}{rrrrrrr}
\toprule
Workers & \multicolumn{3}{c}{small} & \multicolumn{3}{c}{medium} \\
\cmidrule(lr){2-4} \cmidrule(lr){5-7}
 & Time (ms) & $S_{\text{abs}}$ & $f_{\text{parallel}}$ & Time (ms) & $S_{\text{abs}}$ & $f_{\text{parallel}}$ \\
\midrule
1 & 10.0 & 0.93 & --- & 722.7 & 1.01 & --- \\
2 & 8.0 & 1.17 & 0.29 & 683.3 & 1.06 & 0.12 \\
4 & 8.0 & 1.17 & 0.19 & 718.0 & 1.01 & 0.02 \\
8 & 10.3 & 0.90 & -0.12 & 592.7 & 1.23 & 0.21 \\
16 & 12.0 & 0.78 & -0.30 & 592.7 & 1.23 & 0.20 \\
32 & 17.0 & 0.55 & -0.85 & 553.0 & 1.31 & 0.25 \\
64 & 12.0 & 0.78 & -0.29 & 550.7 & 1.32 & 0.25 \\
\bottomrule
\end{tabular}
\end{table*}

\begin{table*}[t]
\caption{OpenMP on maple ($k=3$)}
\label{tab:maple-openmp}
\centering
\begin{tabular}{rrrrrrr}
\toprule
Workers & \multicolumn{3}{c}{small} & \multicolumn{3}{c}{medium} \\
\cmidrule(lr){2-4} \cmidrule(lr){5-7}
 & Time (ms) & $S_{\text{abs}}$ & $f_{\text{parallel}}$ & Time (ms) & $S_{\text{abs}}$ & $f_{\text{parallel}}$ \\
\midrule
1 & 10.3 & 0.90 & --- & 662.0 & 1.10 & --- \\
2 & 15.3 & 0.61 & -1.29 & 577.3 & 1.26 & 0.41 \\
4 & 6.7 & 1.40 & 0.38 & 561.7 & 1.29 & 0.30 \\
8 & 13.7 & 0.68 & -0.53 & 553.3 & 1.31 & 0.27 \\
16 & 10.3 & 0.90 & -0.11 & 565.3 & 1.29 & 0.24 \\
32 & 14.3 & 0.65 & -0.55 & 535.3 & 1.36 & 0.27 \\
64 & 20.0 & 0.47 & -1.16 & 555.0 & 1.31 & 0.24 \\
\bottomrule
\end{tabular}
\end{table*}

\begin{table*}[t]
\caption{MPI on maple ($k=3$)}
\label{tab:maple-mpi}
\centering
\begin{tabular}{rrrrrrr}
\toprule
Workers & \multicolumn{3}{c}{small} & \multicolumn{3}{c}{medium} \\
\cmidrule(lr){2-4} \cmidrule(lr){5-7}
 & Time (ms) & $S_{\text{abs}}$ & $f_{\text{parallel}}$ & Time (ms) & $S_{\text{abs}}$ & $f_{\text{parallel}}$ \\
\midrule
1 & 6.7 & 1.40 & --- & 630.0 & 1.15 & --- \\
2 & 5.0 & 1.87 & 0.93 & 626.3 & 1.16 & 0.28 \\
4 & 12.3 & 0.76 & -0.43 & 758.7 & 0.96 & -0.06 \\
8 & 18.3 & 0.51 & -1.10 & 721.0 & 1.01 & 0.01 \\
16 & 7.0 & 1.33 & 0.27 & 684.3 & 1.06 & 0.06 \\
32 & 29.3 & 0.32 & -2.21 & 932.7 & 0.78 & -0.29 \\
\bottomrule
\end{tabular}
\end{table*}

\begin{figure}[htbp]
\centering
\begin{tikzpicture}
\begin{axis}[
    width=\linewidth, height=5.5cm,
    xlabel={Workers}, ylabel={Absolute speedup},
    xmode=log, log basis x=2,
    legend pos=north west, legend style={font=\scriptsize},
]
\addplot table[x=workers, y=speedup_abs] {data/maple_threaded_medium.dat};
\addplot table[x=workers, y=speedup_abs] {data/maple_openmp_medium.dat};
\addplot table[x=workers, y=speedup_abs] {data/maple_mpi_medium.dat};
\legend{Pthreads, OpenMP, MPI}
\end{axis}
\end{tikzpicture}
\caption{Absolute speedup vs. worker count on maple, medium dataset.}
\label{fig:maple-speedup}
\end{figure}

\subsection{Runtime and Speedup on hickory}

Tables~\ref{tab:hickory-threaded}--\ref{tab:hickory-mpi} report mean
runtime, absolute speedup, and estimated parallel fraction for all three
implementations on hickory, across all three datasets and the worker
counts hickory's 6-CPU job cap allows. Table~\ref{tab:hickory-mpi} reports
only the single-node case, since hickory has no second node to place a
multi-node MPI run on, so the single-node-versus-multi-node comparison
Section~\ref{sec:methodology} describes for Athena does not appear here.
Figure~\ref{fig:hickory-speedup} plots absolute speedup against worker
count for all three implementations on the large dataset.

\begin{table*}[t]
\caption{Pthreads on hickory ($k=3$)}
\label{tab:hickory-threaded}
\centering
\begin{tabular}{rrrrrrrrrr}
\toprule
Workers & \multicolumn{3}{c}{small} & \multicolumn{3}{c}{medium} & \multicolumn{3}{c}{large} \\
\cmidrule(lr){2-4} \cmidrule(lr){5-7} \cmidrule(lr){8-10}
 & Time (ms) & $S_{\text{abs}}$ & $f_{\text{parallel}}$ & Time (ms) & $S_{\text{abs}}$ & $f_{\text{parallel}}$ & Time (ms) & $S_{\text{abs}}$ & $f_{\text{parallel}}$ \\
\midrule
1 & 6.7 & 0.85 & --- & 429.0 & 0.97 & --- & 8286.0 & 0.98 & --- \\
2 & 3.0 & 1.89 & 0.94 & 236.0 & 1.77 & 0.87 & 4177.3 & 1.94 & 0.97 \\
4 & 3.7 & 1.55 & 0.47 & 178.7 & 2.34 & 0.76 & 2204.7 & 3.67 & 0.97 \\
6 & 2.7 & 2.12 & 0.64 & 171.3 & 2.44 & 0.71 & 1537.0 & 5.26 & 0.97 \\
\bottomrule
\end{tabular}
\end{table*}

\begin{table*}[t]
\caption{OpenMP on hickory ($k=3$)}
\label{tab:hickory-openmp}
\centering
\begin{tabular}{rrrrrrrrrr}
\toprule
Workers & \multicolumn{3}{c}{small} & \multicolumn{3}{c}{medium} & \multicolumn{3}{c}{large} \\
\cmidrule(lr){2-4} \cmidrule(lr){5-7} \cmidrule(lr){8-10}
 & Time (ms) & $S_{\text{abs}}$ & $f_{\text{parallel}}$ & Time (ms) & $S_{\text{abs}}$ & $f_{\text{parallel}}$ & Time (ms) & $S_{\text{abs}}$ & $f_{\text{parallel}}$ \\
\midrule
1 & 6.0 & 0.94 & --- & 420.7 & 0.99 & --- & 8142.3 & 0.99 & --- \\
2 & 3.0 & 1.89 & 0.94 & 244.7 & 1.71 & 0.83 & 4122.7 & 1.96 & 0.98 \\
4 & 10.7 & 0.53 & -1.18 & 169.3 & 2.47 & 0.79 & 2112.0 & 3.83 & 0.99 \\
6 & 15.3 & 0.37 & -2.05 & 146.7 & 2.85 & 0.78 & 1638.3 & 4.94 & 0.96 \\
\bottomrule
\end{tabular}
\end{table*}

\begin{table*}[t]
\caption{MPI on hickory, single node ($k=3$)}
\label{tab:hickory-mpi}
\centering
\begin{tabular}{rrrrrrrrrr}
\toprule
Workers & \multicolumn{3}{c}{small} & \multicolumn{3}{c}{medium} & \multicolumn{3}{c}{large} \\
\cmidrule(lr){2-4} \cmidrule(lr){5-7} \cmidrule(lr){8-10}
 & Time (ms) & $S_{\text{abs}}$ & $f_{\text{parallel}}$ & Time (ms) & $S_{\text{abs}}$ & $f_{\text{parallel}}$ & Time (ms) & $S_{\text{abs}}$ & $f_{\text{parallel}}$ \\
\midrule
1 & 5.0 & 1.13 & --- & 418.7 & 1.00 & --- & 8121.7 & 1.00 & --- \\
2 & 2.7 & 2.12 & 1.06 & 212.0 & 1.97 & 0.99 & 4072.3 & 1.99 & 0.99 \\
4 & 4.3 & 1.31 & 0.31 & 134.0 & 3.12 & 0.91 & 2055.0 & 3.94 & 0.99 \\
6 & 3.0 & 1.89 & 0.56 & 92.3 & 4.53 & 0.93 & 1385.7 & 5.84 & 0.99 \\
\bottomrule
\end{tabular}
\end{table*}

\begin{figure}[htbp]
\centering
\begin{tikzpicture}
\begin{axis}[
    width=\linewidth, height=5.5cm,
    xlabel={Workers}, ylabel={Absolute speedup},
    xmode=log, log basis x=2,
    legend pos=north west, legend style={font=\scriptsize},
]
\addplot table[x=workers, y=speedup_abs] {data/hickory_threaded_large.dat};
\addplot table[x=workers, y=speedup_abs] {data/hickory_openmp_large.dat};
\addplot table[x=workers, y=speedup_abs] {data/hickory_mpi_large.dat};
\legend{Pthreads, OpenMP, MPI}
\end{axis}
\end{tikzpicture}
\caption{Absolute speedup vs. worker count on hickory, large dataset.}
\label{fig:hickory-speedup}
\end{figure}

\section{Analysis and Discussion}
\label{sec:analysis}

The two clusters tell opposite stories from the same code, and the
difference is attributable to resource scheduling rather than to any
implementation. On hickory, every job receives a dedicated Slurm
allocation that no other job touches, and
Tables~\ref{tab:hickory-threaded}--\ref{tab:hickory-mpi} show the
consequence directly: all three
implementations scale consistently with worker count on the medium and
large datasets, reaching between 4.9$\times$ and 5.8$\times$ speedup at 6
workers on the large dataset, with an estimated parallel fraction above
0.9 for nearly every medium and large configuration. That high parallel
fraction matches the algorithm's structure, since the only serial work in
any implementation is the single confusion-matrix pass over the finished
predictions, an $O(\mathit{num\_classes}^2)$ operation that is negligible
next to the $O(\mathit{num\_test} \times \mathit{num\_train} \times
\mathit{num\_attributes})$ distance computation every worker performs in
parallel.

maple's shared, unscheduled hardware shows a materially different picture
on the same algorithm. On the medium dataset, all three implementations
still show real but far more modest gains, topping out around
1.3$\times$--1.4$\times$ speedup for pthreads and OpenMP, well short of
hickory's numbers at similar worker counts. MPI on maple's medium dataset
is worse still: speedup is roughly flat through 16 workers and then drops
below 1 at 32 workers (Table~\ref{tab:maple-mpi}), meaning 32 concurrent
MPI ranks finished slower than one. MPI ranks are heavier OS entities than
threads, each with its own scheduling footprint and its own participation
in the \texttt{MPI\_Barrier} and \texttt{MPI\_Gatherv} synchronization
points described in Section~\ref{sec:methodology}; on a machine already
shared with other users' processes, oversubscribing it with 32 such ranks
multiplies contention for the same physical cores far more than 32 threads
inside one process would, and every synchronization point then waits on
whichever rank the scheduler starved the longest.

On the small dataset this effect is severe enough to produce negative
estimated parallel fractions in every maple table, for example OpenMP at
64 workers (Table~\ref{tab:maple-openmp}) taking 20.0\,ms against a 9.3\,ms
serial baseline, slower with 64 workers than with one. Small's runs
complete in only a few milliseconds, so the fixed cost of creating threads
or ranks and competing with other users' processes for scarce cores
dominates a workload that is barely there to begin with; hickory shows
traces of the same effect at its own highest worker count on small
(Tables~\ref{tab:hickory-threaded}--\ref{tab:hickory-openmp}), just far
less severely, since a dedicated allocation still removes the
cross-user-contention half of the problem. Karp-Flatt is doing its job
correctly in both cases: a negative estimated parallel fraction is not a
nonsensical result but an honest signal that overhead, not serial code,
is what is limiting speedup at that configuration.

Comparing the three implementations directly on hickory's large dataset at
6 workers, the highest worker count tested, MPI reaches 5.84$\times$,
pthreads 5.26$\times$, and OpenMP 4.94$\times$. The three are within
roughly 15\% of each other, consistent with all three performing the same
per-worker distance computation once started; the gap is not large enough,
relative to the three repetitions each figure is already averaged over, to
claim a definitive ranking between them, only that none of the three pays
a heavy implementation-specific tax on hickory's dedicated hardware.

Two comparisons the assignment specifies could not be made in this report.
The single-node-versus-multi-node MPI comparison has no hickory
counterpart, since hickory exposes only one node; Table~\ref{tab:hickory-mpi}
reports the single-node case only, and no statement can be made here about
inter-node communication cost. Athena itself, which would have supported
that comparison, was inaccessible for the entire assignment period despite
being reported to the course staff, so neither cluster's data in this
report can speak to it. Separately, the worker-count ranges tested differ
by cluster for two different reasons: hickory's ceiling of 6 is an
administrative QOS limit, not a technical one, while maple's practical
ceiling on the two larger datasets is the oversubscription effect
described above and in Section~\ref{sec:methodology}. Both are themselves
findings about how scheduling policy and resource contention, not the
parallelization code, bound what is achievable in practice.

\section{Conclusion}
\label{sec:conclusion}

This report implemented and evaluated three parallel strategies, POSIX
threads, OpenMP, and MPI, for the k-nearest-neighbor classifier, all built
around the same test-instance partitioning strategy. That shared strategy
was chosen specifically because it reproduces the serial version's
predictions exactly regardless of worker count, which was confirmed
empirically across all three implementations, three datasets, and all 244
completed runs on both clusters used in this study.

The results show that the same parallelization code can succeed or fail
depending entirely on the hardware it runs on. On hickory's dedicated,
uncontended Slurm allocation, all three implementations scaled
consistently with worker count, reaching up to 5.8$\times$ speedup on the
large dataset at 6 workers with an estimated parallel fraction above 0.9,
matching KNN's embarrassingly parallel structure. On maple's shared,
unscheduled hardware, the same code produced only modest gains on the
medium dataset and measurable slowdowns on the small one, since the fixed
cost of creating threads or MPI ranks, and competing with other users'
processes for the same physical cores, outweighed the benefit of
parallelizing workloads that were themselves only milliseconds long.
Athena, which the assignment specifies for the single-node-versus-multi-node
MPI comparison, was inaccessible for the entire assignment period despite
being reported to the course staff; the hickory results substituted here
are otherwise complete but, having only one physical node, cannot speak to
that specific comparison.

\begin{thebibliography}{9}
\bibitem{knn} T. Cover and P. Hart, ``Nearest neighbor pattern classification,'' \emph{IEEE Transactions on Information Theory}, vol. 13, no. 1, pp. 21--27, 1967.
\bibitem{pthreads} B. Lewis and D. J. Berg, \emph{Multithreaded Programming with Pthreads}. Prentice-Hall, Inc., 1998.
\bibitem{openmp} OpenMP Architecture Review Board, ``OpenMP Application Programming Interface, Version 5.2,'' 2021.
\bibitem{mpi} W. Gropp, E. Lusk, and A. Skjellum, \emph{Using MPI: Portable Parallel Programming with the Message-Passing Interface}, 3rd ed. MIT Press, 2014.
\bibitem{karpflatt} A. H. Karp and H. P. Flatt, ``Measuring parallel processor performance,'' \emph{Communications of the ACM}, vol. 33, no. 5, pp. 539--543, 1990.
\end{thebibliography}

\end{document}
